\section{Preliminary}\label{sec:Preliminary}
    % 写预备知识，主要是别人的论文里面的已经提出的方法、技术、或者损失，在本文中要用到。但是明确不是自己提出的。
    % 写一段简短的开场白。注意，这里开始，给每个东西所使用的名字、数学符号，就得定好，后续方法章节也得统一符号明明。最好给出一个表格，实时记录每个变量的含义。
    \subsection{Notations}\label{sec:Preliminary:Notations}
        To avoid ambiguity, the main symbols used in this paper are summarized in Table~\ref{tab:notations}. Their definitions are fixed throughout the remainder of the manuscript.

        \begin{table}[t]
            \centering
            \caption{Summary of the main symbols used in this paper.}
            \label{tab:notations}
            \footnotesize
            \setlength{\tabcolsep}{4pt}
            \renewcommand{\arraystretch}{1.15}
            \begin{tabular}{>{\raggedright\arraybackslash}p{0.28\columnwidth}>{\raggedright\arraybackslash}p{0.62\columnwidth}}
                \hline
                Symbol & Meaning \\
                \hline
                $\mathbf{I}^{t_1}, \mathbf{I}^{t_2}$ & Bi-temporal input images captured at times $t_1$ and $t_2$. \\
                $\hat{\mathbf{Y}}$ & Predicted binary change map. \\
                $\mathbf{Y}$ & Ground-truth binary change map. \\
                $\mathbf{Y}^{e}$ & Ground-truth edge supervision map. \\
                $\mathbf{P}\in\mathbb{R}^{2\times H\times W}$ & Change prediction logits. \\
                $\mathbf{P}^{e}\in\mathbb{R}^{2\times H\times W}$ & Edge prediction logits. \\
                $\mathbf{F}^{t_1}, \mathbf{F}^{t_2}\in\mathbb{R}^{C\times H\times W}$ & Deep feature maps extracted from the two temporal observations. \\
                $\mathcal{N}_{\theta}$ & Change detection network parameterized by $\theta$. \\
                $\Omega=\{1,\ldots,H\}\times\{1,\ldots,W\}$ & Spatial domain of an image. \\
                $(i,j)\in\Omega$ & Pixel location at row $i$ and column $j$. \\
                $p^{1}_{b,ij}, q^{1}_{b,ij}$ & Foreground probabilities of the change and edge branches after softmax. \\
                $B$ & Batch size. \\
                $\alpha$ & Balance factor between CE loss and IoU loss within the segmentation branch. \\
                $\beta$ & Weight of the contrastive branch in the final objective. \\
                $\gamma$ & Weight of the edge branch in the final objective. \\
                $m$ & Contrastive margin. \\
                $\varepsilon$ & Numerical stabilizer. \\
                \begin{tabular}[t]{@{}l@{}}
                $\mathcal{L}_{\text{ce}}, \mathcal{L}_{\text{iou}}, \mathcal{L}_{\text{edge}}$ \\
                $\mathcal{L}_{\text{con}}, \mathcal{L}_{\text{seg}}, \mathcal{L}_{\text{total}}$
                \end{tabular}
                & Cross-entropy, IoU, edge, contrastive, segmentation-fusion, and total optimization losses. \\
                \hline
            \end{tabular}
        \end{table}
    
    \subsection{Region Changing Detection}\label{sec:Preliminary:Region Changing Detection}
        % 用公式和段落，介绍遥感领域汇总的伪变化检测的定义。
        Remote sensing change detection aims to identify whether each spatial location has undergone semantic change between two observations. Given a bi-temporal image pair $\left(\mathbf{I}^{t_1}, \mathbf{I}^{t_2}\right)$, the task is to predict a binary change map $\hat{\mathbf{Y}}\in\{0,1\}^{H\times W}$, where $0$ denotes unchanged and $1$ denotes changed. In practical benchmark datasets, the target is not merely to respond to visual discrepancy, but to produce predictions consistent with the annotation criteria defined by the dataset.
    
    \subsection{Pseudo-Change}\label{sec:Preliminary:Pseudo-Change}
        % 用公式以及段落，介绍伪变化检测的定义。
        In remote sensing scenes, many regions exhibit visible appearance differences between two acquisition times even though the semantic land-cover category remains unchanged. These differences may be caused by illumination variation, seasonal transition, atmospheric interference, sensor difference, or local imaging conditions. When such regions are still annotated as unchanged in the ground truth, they form a pseudo-change phenomenon.

        Formally, for a pixel location $(i,j)$, let $\Delta v_{ij}$ denote the visual discrepancy between the bi-temporal observations and let $y_{ij}\in\{0,1\}$ denote the benchmark annotation. A pseudo-change region satisfies

        \begin{equation}
            \Delta v_{ij} > 0, \qquad y_{ij} = 0,
        \end{equation}
        which means that observable appearance variation exists, but the annotation semantics still indicate unchanged status.

        Pseudo-change is one of the major causes of false positives in change detection. If a model is optimized only to amplify visual differences, it may over-respond to non-semantic variations and generate change predictions inconsistent with benchmark labels. Therefore, suppressing pseudo-change interference is essential for improving annotation-consistent change detection.

    \subsection{Segmentation Loss Computation}\label{sec:Preliminary:SegmentationLossComputation}
        This work employs several standard supervised losses as basic optimization components. These losses are introduced here as preliminary definitions, whereas the proposed contribution lies in their bounded normalization and coupling with feature-level contrastive supervision in the Methodology section. Let $\Omega=\{1,\ldots,H\}\times\{1,\ldots,W\}$ denote the spatial domain of an image, and let $(i,j)\in\Omega$ denote the pixel at row $i$ and column $j$. Here, $\mathbf{P}$ and $\mathbf{P}^{e}$ denote the logits of the change and edge branches, respectively, while $\mathbf{Y}$ and $\mathbf{Y}^{e}$ denote the corresponding binary supervision maps. For a mini-batch of size $B$, the foreground probabilities of the change and edge branches are given by
        \begin{equation}
            \begin{aligned}
            p^{1}_{b,ij}
            &=
            \frac{\exp\left(\mathbf{P}_{b,1,i,j}\right)}
            {\sum_{c=0}^{1}\exp\left(\mathbf{P}_{b,c,i,j}\right)}, \\
            q^{1}_{b,ij}
            &=
            \frac{\exp\left(\mathbf{P}^{e}_{b,1,i,j}\right)}
            {\sum_{c=0}^{1}\exp\left(\mathbf{P}^{e}_{b,c,i,j}\right)}.
            \end{aligned}
            \label{eq:foreground_probability}
        \end{equation}

        For compactness, $\sum_{b,ij}$ denotes $\sum_{b=1}^{B}\sum_{(i,j)\in\Omega}$ in the following definitions. The cross-entropy loss for pixel-wise change classification is defined as
        \begin{equation}
            \mathcal{L}_{\text{ce}}
            =
            -\frac{1}{B|\Omega|}
            \sum_{b,ij}
            \log
            \frac{\exp\left(\mathbf{P}_{b,y^{b}_{ij},i,j}\right)}
            {\sum_{c=0}^{1}\exp\left(\mathbf{P}_{b,c,i,j}\right)}.
            \label{eq:ce_loss}
        \end{equation}

        The foreground IoU loss for region-level overlap supervision is defined as
        \begin{equation}
            \mathcal{L}_{\text{iou}}
            =
            1-\frac{I_{\text{chg}}}{U_{\text{chg}}+\varepsilon},
            \label{eq:iou_loss}
        \end{equation}
        where
        \begin{equation}
            I_{\text{chg}}=\sum_{b,ij}p^{1}_{b,ij}y^{b}_{ij},
            \qquad
            U_{\text{chg}}=\sum_{b,ij}\left(p^{1}_{b,ij}+y^{b}_{ij}-p^{1}_{b,ij}y^{b}_{ij}\right).
            \label{eq:iou_components}
        \end{equation}

        The foreground Dice loss for edge-aware supervision is defined as
        \begin{equation}
            \mathcal{L}_{\text{edge}}
            =
            1-\frac{2I_{\text{edge}}+\varepsilon}{S_{\text{edge}}+\varepsilon},
            \label{eq:edge_loss}
        \end{equation}
        where
        \begin{equation}
            I_{\text{edge}}=\sum_{b,ij}q^{1}_{b,ij}y^{e,b}_{ij},
            \qquad
            S_{\text{edge}}=\sum_{b,ij}\left(q^{1}_{b,ij}+y^{e,b}_{ij}\right).
            \label{eq:edge_components}
        \end{equation}
        Here, $\varepsilon=10^{-7}$ is a small constant for numerical stability. Since the implementation uses \texttt{num\_classes=2}, the IoU and Dice losses are computed only on the foreground class.